\documentclass[12pt]{article}
%---DOCUMENT MARGINS---
\usepackage{geometry} % Required for adjusting page dimensions and margins
\geometry{
	paper=a4paper, % Paper size, change to letterpaper for US letter size
	top=4cm, % Top margin
	bottom=2cm, % Bottom margin
	left=2cm, % Left margin
	right=2cm, % Right margin
	headheight=3cm, % Header height
	%footskip=1.5cm, % Space from the bottom margin to the baseline of the footer
	headsep=0.5cm, % Space from the top margin to the baseline of the header
	%showframe, % Uncomment to show how the type block is set on the page
}
\usepackage{cmap}

\usepackage[T1,T2A]{fontenc}
\usepackage{fontspec}
\setmainfont[
	BoldFont={* Bold},
	ItalicFont={* Italic},
	BoldItalicFont={* BoldItalic}
]{DejaVu Serif Condensed}
\setsansfont{DejaVu Sans}
\setmonofont{Dejavu Sans Mono}
\usepackage[math-style=TeX]{unicode-math}
\setmathfont{Latin Modern Math}

\usepackage[nobottomtitles*]{titlesec}
\titleformat
{\section} % command
{\fontsize{14}{14}\sffamily\bfseries} % format
{} % label
{0pt} % sep
{} % before-code
\titlespacing{\section}{0pt}{0.75em}{0.25em}
\newcommand{\tl}{}
\newcommand{\ml}{}
\newcommand{\problem}[1]{%
	\section[#1]{#1 \fontsize{14}{14}\selectfont{\hfill{\normalfont{
				\emoji{hourglass-not-done}\tl \space\space \emoji{floppy-disk} \ml}}}}
}
\AddToHook{cmd/section/before}{\clearpage}

\usepackage[dvipsnames]{xcolor}
\titleformat
{\subsection} % command
{\fontsize{14}{14}\sffamily\bfseries} % format
{\includesvg[width=0.035\textwidth, inkscapelatex=false]{lion}\space} % label
{0pt} % sep
{} % before-code
\titlespacing{\subsection}{0pt}{0.75em}{0.25em}

\setlength{\parskip}{0.5em}
\setlength{\parindent}{0pt}
\renewcommand{\baselinestretch}{1.2}
\sloppy

\usepackage{listings}
\definecolor{lstbg}{RGB}{250,250,252}
\definecolor{lstframe}{RGB}{220,220,225}
\lstdefinestyle{mystyle}{
	backgroundcolor=\color{lstbg},
	frame=single,
	rulecolor=\color{lstframe},
	basicstyle=\ttfamily\small,
	keywordstyle=\bfseries,
	breaklines=true,
	aboveskip=0.5\baselineskip,
	belowskip=0\baselineskip  
}
\lstset{style=mystyle, language=C++}

\usepackage{fancyhdr}
\pagestyle{fancy}
\setlength{\headwidth}{\textwidth}
\newcommand{\round}{}
\newcommand{\problemName}{}
\newcommand{\lang}{English}
\fancyhead[L]{
	\begin{minipage}{0.4\textwidth}
		\bf{
			EJOI 2025 \round\\
			\problemName\\
			\emoji{globe-with-meridians} \lang
		}
	\end{minipage}
	\vspace{0.5cm}
}
\fancyhead[R]{%
	\begin{minipage}{0.6\textwidth}
		\includesvg[width=\textwidth, inkscapelatex=false]{logo}
	\end{minipage}
	\vspace{0.5cm}
}
\usepackage{lastpage}
\fancyfoot[C]{\problemName\space(\lang) \thepage\ / \pageref*{LastPage}}

\raggedbottom

\usepackage{amsmath}
\usepackage{stmaryrd}

\usepackage{graphicx}
\graphicspath{{./}}
\usepackage[export]{adjustbox}
\usepackage{wrapfig}
\makeatletter
\patchcmd\WF@putfigmaybe{\lower\intextsep}{}{}{\fail}
\AddToHook{env/wrapfigure/begin}{\setlength{\intextsep}{0pt}}
\makeatother
\usepackage[inkscapearea=page, inkscapepath=./svg-inkscape]{svg}
\svgpath{{./}}

\usepackage{makecell}
\usepackage{tabularray}
\AtBeginEnvironment{table}{\vspace{-0.2cm}}
\AtEndEnvironment{table}{\vspace{-0.2cm}}
\usepackage{float}
\usepackage{placeins}
\usepackage{caption}
\captionsetup[table]{
	skip=0.25em,
	singlelinecheck=false, justification=justified
}

\usepackage{enumitem}
\setlist{itemsep=-0.4em, leftmargin=22pt, topsep=-\parskip}
\AfterEndEnvironment{itemize}{\vspace{\parskip}}
\newcommand{\tabitem}{\indent~~\llap{\textbullet}~~}

\usepackage{hyperref}
\hypersetup{
	colorlinks=true,
	citecolor=blue,
	linkcolor=blue,
	urlcolor=cyan,
}

\usepackage{emoji}

\usepackage[english]{babel}

\begin{document}

\renewcommand{\round}{Day 1}
\renewcommand{\problemName}{Task Prison}
\renewcommand{\lang}{Russian}

\renewcommand{\tl}{$2$ sec.}
\renewcommand{\ml}{$1024$ MB}
\problem{\problemName}

Алису и Боба несправедливо приговорили к тюремному заключению строгого режима. Теперь им предстоит спланировать побег. Для этого им необходимо общаться максимально эффективно (в частности, Алисе необходимо ежедневно отправлять Бобу информацию). Однако они не могут встречаться и могут обмениваться информацией только посредством записей, написанных на салфетках. Каждый день Алиса хочет отправлять Бобу новую информацию — число от $0$ до $N-1$. За каждым обедом Алиса берёт три салфетки, пишет на каждой из них число от $0$ до $M-1$ (могут быть повторения) и оставляет их на своём стуле. Затем их враг, Чарли, уничтожает одну из салфеток и перемешивает две другие. Наконец, Боб находит две оставшиеся салфетки и читает на них числа. Он должен точно расшифровать исходное число, которое Алиса хотела ему отправить. Место на салфетках ограничено, поэтому $M$ фиксировано. Однако цель Алисы и Боба — максимизировать пропускную способность, поэтому они могут выбрать $N$ как можно большего значения. Помогите Алисе и Бобу, реализовать стратегию для каждого из них, чтобы максимизировать значение $N$.

\subsection{Implementation details}
Поскольку это задача коммуникации, ваша программа будет запущена в двух отдельных запусках (одно для Алисы и одно для Боба), которые не смогут обмениваться данными или взаимодействовать каким-либо другим способом, кроме описанного здесь. Вам необходимо реализовать три функции:

\begin{lstlisting}
 int setup(int M);
\end{lstlisting}

Эта функция будет вызвана один раз при запуске вашей программы Алисой и один раз при запуске программы Бобом. Функция получает $M$ и должна вернуть желаемое $N$. Оба вызова \texttt{setup} должны возвращать одинаковое $N$.

\begin{lstlisting}
 std::vector<int> encode(int A);
\end{lstlisting}

Это реализует стратегию Алисы. Функция будет вызвана с числом, закодированным в $A$ ($0 \leq A < N$), и должна вернуть три числа $W_1, W_2, W_3$ ($0 \leq W_i < M$), закодировавших $A$. Эта функция будет вызвана всего $T$ раз — один раз в день (значения $A$ могут повторяться в разные дни).

\begin{lstlisting}
 int decode(int X, int Y);
\end{lstlisting}

Это реализует стратегию Боба. Функция будет вызвана с двумя из трёх чисел, возвращаемых функцией \texttt{encode}, в определённом порядке. Она должна вернуть то же значение $A$, что и функция \texttt{encode}. Эта функция также будет вызвана $T$  раз — что соответствует $T$ вызовам функции \texttt{encode}; они будут выполнены в том же порядке. Все вызовы \texttt{encode} будут выполняться перед вызовами \texttt{decode}.

\subsection{Constraints}

\begin{itemize}
\item $M \leq 4300$
\item $T = 5000$
\end{itemize}

\subsection{Scoring}

Для конкретной подзадачи доля S полученных баллов зависит от наименьшего $N$, возвращаемого \texttt{setup} для любого теста в этой подзадаче. Она также зависит от $N^*$ — целевого значения $N$, необходимого для получения полного количества баллов за подзадачу:

\begin{itemize}
\item Если ваше решение не проходит хотя бы один тест, то $S = 0$.
\item Если $N \geq N^*$, то $S = 1,0$.
\item Если $N < N^*$, то $S = \max\left(0.35 \max\left(\frac{\log(N) - 0.985 \log(M)}{\log(N^*) - 0.985 \log(M)}, 0.0\right)^{0.3} + 0.65 \left(\frac{N}{N^*}\right)^{2.4}, 0.01\right)$.
\end{itemize}

\subsection{Subtasks}

\begin{table}[H]
\begin{tblr}{|Q[c,m]|Q[c,m]|Q[c,m]|Q[c,m]|}
\hline
\textbf{Подзадача} & \textbf{Баллы} & $M$ & $N^*$ \\
\hline
$1$ & $10$ & $700$ & $82017$ \\
\hline
$2$ & $10$ & $1100$ & $202217$ \\
\hline
$3$ & $10$ & $1500$ & $375751$ \\
\hline
$4$ & $10$ & $1900$ & $602617$ \\
\hline
$5$ & $10$ & $2300$ & $882817$ \\
\hline
$6$ & $10$ & $2700$ & $1216351$ \\
\hline
$7$ & $10$ & $3100$ & $1603217$ \\
\hline
$8$ & $10$ & $3500$ & $2043417$ \\
\hline
$9$ & $10$ & $3900$ & $2536951$ \\
\hline
$10$ & $10$ & $4300$ & $3083817$ \\
\hline
\end{tblr}
\end{table}
\FloatBarrier

\subsection{Example}

Рассмотрим следующий пример с $T = 5$. Здесь мы имеем схему кодирования, в которой Алиса отправляет три одинаковых числа для кодирования $0$ или три различных числа для кодирования $1$. Обратите внимание, что Боб может декодировать исходное число из любых двух из трёх чисел, отправленных Алисой.

\begin{table}[H]
\begin{tblr}{|Q[c,m]|Q[c,m]|Q[c,m]|}
\hline \textbf{\makecell{Запуск}} & \textbf{\makecell{Вызов функции}} & \textbf{Возвращаемое значение} \\
\hline
Alice & \texttt{setup(10)} & \texttt{2} \\
\hline
Bob & \texttt{setup(10)} & \texttt{2} \\
\hline
Alice & \texttt{encode(0)} & \texttt{\{5, 5, 5\}} \\
\hline
Alice & \texttt{encode(1)} & \texttt{\{8, 3, 7\}} \\
\hline
Alice & \texttt{encode(1)} & \texttt{\{0, 3, 1\}} \\
\hline
Alice & \texttt{encode(0)} & \texttt{\{7, 7, 7\}} \\
\hline
Alice & \texttt{encode(1)} & \texttt{\{6, 2, 0\}} \\
\hline
Bob & \texttt{decode(5, 5)} & \texttt{0} \\
\hline
Bob & \texttt{decode(8, 7)} & \texttt{1} \\
\hline
Bob & \texttt{decode(3, 0)} & \texttt{1} \\
\hline
Bob & \texttt{decode(7, 7)} & \texttt{0} \\
\hline
Bob & \texttt{decode(2, 0)} & \texttt{1} \\
\hline
\end{tblr}
\end{table}

\subsection{Sample grader}

В примере оценщика все вызовы \texttt{encode} и \texttt{decode} будут происходить в одном и том же выполнении программы. Кроме того, \texttt{setup} будет вызван только один раз (а не дважды, по одному разу за выполнение, как в системе оценивания).

На вход поступает одно целое число — $M$. Затем программа выведет $N$, которые вы вернули в вашей программе \texttt{setup}. Затем она вызовет функции \texttt{encode} и \texttt{decode} в указанном порядке $T$ раз со случайно сгенерированными числами от $0$ до $N-1$ и случайно сгенерированными вариантами выбора двух из трёх чисел из \texttt{encode} для функции \texttt{decode} (и в каком порядке). Если решение не удалось, программа выведет сообщение об ошибке.

\end{document}
